% CHAPTER 1: INTRODUCTION
% ======================================================================
\setstretch{1.5}
\pagenumbering{arabic}
\setcounter{page}{1}
{\fontsize{18}{22}\selectfont\textbf{1 Introduction}}\par

\vspace{3mm}

\subhead{1.1 Background and Motivation}

LLMs are built on transformer architectures. Through training on large sets of text, these systems can understand almost every language, create runnable code from a text prompt, and give detailed answers about many fields. This is why engineers see them as a key piece of enterprise digital change.

But using LLMs in the banking industry brings several technical and non-technical problems. First, LLM-generated SQL is not always reliable. The code often has mistakes in structure or meaning, especially when the database has a complex layout. These problems may lead to incorrect changes to funds. Second, AI models have limited knowledge of bank-specific rules. Frequent updates to regulations and policies make the model's knowledge outdated. Third, bank workers must integrate AI into their daily work, such as room booking, schedule checking, and route planning. However, present AI models are not always good at calling the right tool or reading the right input. All these problems can lead to operational risk unless the AI system is constrained by security, retrieval, and tool-execution controls.

The current work is driven by several linked problems. These are technical precision, keeping domain knowledge correct, stable tool calling, and data control. So, a single multi-agent AI system is clearly needed. This system must work on-site and take a defense-in-depth approach. It must combine secure code generation, local knowledge retrieval, and enterprise tool calling.

\subhead{1.2 Problem Statement}

The present work is facing four major challenges.

\textbf{Code generation accuracy.} Database queries are used in everyday bank work. For example, they are used to get customer data and to file reports. Natural-language interfaces can make data access much easier. But LLM-written SQL is not always safe. It must be run through careful checks. The model may add syntax faults, wrong column references, or even harmful operations. The system needs to check multiple levels before executing any SQL statement.

\textbf{Knowledge retrieval.} Answers to procedural and policy questions must be verifiable. This is important because finance is a field where rules must be followed and there is zero tolerance for errors. RAG methods are used to fix LLM issues with old or made-up information. But bank documents have security levels. The system must filter retrieval results to match the user's clearance level. Documents must not leak through the agent's answers.

\textbf{Enterprise Tool Integration.} Daily bank work uses many tools, including room booking, calendar syncing, and travel route planning. An AI system must correctly understand user requests. It must find the right intent, extract the right parameters, and run the matching tool. The system should also use LLM routing to pick the best handler for each request.

\textbf{Data Sovereignty And Access Control.} Beyond the risks from outside APIs, banking settings need strict access levels. Different departments handle documents with different security clearances. The rule of least privilege must be enforced at every layer. Small permission models are often too simple for an organization that has many departments and levels of leadership.

\subhead{1.3 Project Objectives}

This work presents BankAgent, a multi-agent AI platform built on on-premise infrastructure and designed for banking operations. The platform delivers three AI agents under a unified Copilot dispatcher, supported by five platform services and one-command containerized deployment. Table \ref{tab:objectives} summarizes the project objectives.

\begin{table}[H]
\centering
\small\singlespacing
\caption{Project Objectives by Platform Layer}
\label{tab:objectives}
\begin{tabularx}{\textwidth}{>{\RaggedRight\arraybackslash}p{0.18\textwidth}>{\RaggedRight\arraybackslash}p{0.28\textwidth}>{\RaggedRight\arraybackslash}X}
\toprule
\textbf{Category} & \textbf{Component} & \textbf{Objective} \\
\midrule
\multirow{3}{*}{AI Agents} & Code Agent & Safely translate natural language questions into executable SQL queries. \\
\cmidrule{2-3}
 & RAG Agent & Retrieve policy knowledge from Milvus-indexed documents with verifiable citations and clearance based permission filtering. \\
\cmidrule{2-3}
 & Tool Agent & Automate business tasks such as room booking, schedule conflict detection, and route planning through natural language. \\
\addlinespace
Intelligent Dispatch & Copilot & Classify user intent from a single interface and dispatch requests to the appropriate agent. \\
\addlinespace
\multirow{5}{*}{Platform Services} & Auth \& RBAC & Enforce JWT authentication with three role tiers and department scoped data access. \\
\cmidrule{2-3}
 & API Gateway & Provide a unified entry point with JWT filtering, rate limiting, and session enforcement. \\
\cmidrule{2-3}
 & Notification System & Manage multi type messages with a five state approval workflow for sensitive operations. \\
\cmidrule{2-3}
 & Document Management & Store and serve documents with tiered security clearance and HITL escalation. \\
\cmidrule{2-3}
 & Task Center & Orchestrate asynchronous agent execution with retry logic and real time progress streaming. \\
\addlinespace
DevOps & Containerized Deploy & Start all services with a single Docker Compose command. \\
\bottomrule
\end{tabularx}
\end{table}

\subhead{1.4 Scope and Limitations}

This project covers the design, implementation, and evaluation of the multi-agent platform described above. The intended users are banking personnel. The platform provides three main functions. It generates and checks SQL queries. It does knowledge retrieval with citations. It calls business tools. The deployment model is local. It uses Docker Compose for container management. The security model includes JWT authentication and RBAC.

The completed system also has defined scope boundaries. First, the Code Agent supports MySQL SQL and relational table structures rather than other database systems or NoSQL interfaces. Second, the RAG Agent's retrieval quality depends on the selected embedding model and the coverage of indexed documents, so the knowledge base must be curated. Third, the Tool Agent supports the three enterprise tools implemented in this project. Fourth, the platform is delivered as a single-instance Docker Compose deployment rather than a Kubernetes or multi-node cluster.

\subhead{1.5 Report Organization}

The rest of this report is arranged like this. Chapter 2 looks at past work on LLM use in business settings, systems with many agents, turning text into SQL, RAG methods, language models with added tools, and on-site use and data safety. Chapter 3 explains what banks need, functional and non-functional requirements, and user case studies. Chapter 4 shows the five-layer system design, the frontend and backend layouts, data layer construction, and the reasons for choosing each technology. Chapter 5 explains in detail how the three main agents are built. Chapter 6 covers the process of putting the system parts together, including user management, task planning, live communication, and deploying with containers. Chapter 7 talks about the project's difficulties, how this work compares to other solutions, and where it falls short. Chapter 8 sums up the completed work and possible extensions beyond the delivered system.

\newpage

% ======================================================================
